\textbf{Phần 1. Trắc nghiệm lựa chọn}
\setcounterpageref{ex}{0}

\begin{ex} %[4D10N57]
Quy đồng mẫu số hai phân số là
\choice
{Cộng mẫu số của hai phân số lại}
{Đổi chỗ hai mẫu số}
{Lấy tử số cộng mẫu số}
{\True Biến đổi hai phân số thành hai phân số mới có cùng mẫu số và lần lượt bằng các phân số đã cho}
\loigiai{
Quy đồng là đưa các phân số về cùng mẫu số.\\
\textbf{Đáp án của câu hỏi là D}
}
\end{ex}

\begin{ex} %[4D10N57]
Mẫu số chung của hai phân số $\dfrac{1}{3}$ và $\dfrac{2}{5}$ là
\choice
{\True $15$}
{$8$}
{$12$}
{$10$}
\loigiai{
$3 \times 5 = 15$.\\
\textbf{Đáp án của câu hỏi là A}
}
\end{ex}

\begin{ex} %[4D10H57]
Quy đồng mẫu số hai phân số $\dfrac{2}{3}$ và $\dfrac{3}{4}$ với mẫu số chung là 12 ta được
\choice
{$\dfrac{6}{12}$ và $\dfrac{9}{12}$}
{\True $\dfrac{8}{12}$ và $\dfrac{9}{12}$}
{$\dfrac{8}{12}$ và $\dfrac{6}{12}$}
{$\dfrac{4}{12}$ và $\dfrac{9}{12}$}
\loigiai{
$\dfrac{2 \times 4}{3 \times 4} = \dfrac{8}{12}$ và $\dfrac{3 \times 3}{4 \times 3} = \dfrac{9}{12}$.\\
\textbf{Đáp án của câu hỏi là B}
}
\end{ex}

\begin{ex} %[4D10H57]
Khi quy đồng mẫu số hai phân số $\dfrac{5}{6}$ và $\dfrac{7}{12}$, mẫu số chung nhỏ nhất thuận tiện là
\choice
{$72$}
{$18$}
{\True $12$}
{$24$}
\loigiai{
Vì 12 chia hết cho 6 nên chọn 12 làm MSC.\\
\textbf{Đáp án của câu hỏi là C}
}
\end{ex}

\begin{ex} %[4D10H57]
Quy đồng mẫu số hai phân số $\dfrac{3}{5}$ và $\dfrac{4}{7}$ ta được hai phân số là
\choice
{$\dfrac{21}{35}$ và $\dfrac{15}{35}$}
{$\dfrac{15}{35}$ và $\dfrac{20}{35}$}
{$\dfrac{21}{35}$ và $\dfrac{28}{35}$}
{\True $\dfrac{21}{35}$ và $\dfrac{20}{35}$}
\loigiai{
$\dfrac{3 \times 7}{5 \times 7} = \dfrac{21}{35}$ và $\dfrac{4 \times 5}{7 \times 5} = \dfrac{20}{35}$.\\
\textbf{Đáp án của câu hỏi là D}
}
\end{ex}

\begin{ex} %[4D10H57]
Mẫu số chung của hai phân số $\dfrac{1}{4}$ và $\dfrac{5}{8}$ là
\choice
{$12$}
{\True $8$}
{$32$}
{$16$}
\loigiai{
8 chia hết cho 4 nên MSC là 8.\\
\textbf{Đáp án của câu hỏi là B}
}
\end{ex}

\begin{ex} %[4D10V57]
Quy đồng mẫu số ba phân số $\dfrac{1}{2}, \dfrac{2}{3}, \dfrac{3}{4}$ với mẫu số chung bằng 12 ta được
\choice
{$\dfrac{6}{12}, \dfrac{4}{12}, \dfrac{9}{12}$}
{$\dfrac{3}{12}, \dfrac{8}{12}, \dfrac{9}{12}$}
{\True $\dfrac{6}{12}, \dfrac{8}{12}, \dfrac{9}{12}$}
{$\dfrac{6}{12}, \dfrac{8}{12}, \dfrac{6}{12}$}
\loigiai{
$\dfrac{6}{12}, \dfrac{8}{12}, \dfrac{9}{12}$.\\
\textbf{Đáp án của câu hỏi là C}
}
\end{ex}

\begin{ex} %[4D10V57]
Quy đồng mẫu số hai phân số $\dfrac{7}{10}$ và $\dfrac{4}{15}$ với mẫu số chung nhỏ nhất là 30 ta được
\choice
{\True $\dfrac{21}{30}$ và $\dfrac{8}{30}$}
{$\dfrac{14}{30}$ và $\dfrac{8}{30}$}
{$\dfrac{21}{30}$ và $\dfrac{12}{30}$}
{$\dfrac{70}{30}$ và $\dfrac{40}{30}$}
\loigiai{
$\dfrac{7 \times 3}{10 \times 3} = \dfrac{21}{30}$ và $\dfrac{4 \times 2}{15 \times 2} = \dfrac{8}{30}$.\\
\textbf{Đáp án của câu hỏi là A}
}
\end{ex}

\begin{ex} %[4D10V57]
Cho hai phân số $\dfrac{5}{8}$ và $\dfrac{3}{4}$. Sau khi quy đồng mẫu số chung là 8 thì tổng hai tử số mới là
\choice
{$10$}
{$12$}
{\True $11$}
{$13$}
\loigiai{
$\dfrac{5}{8}$ và $\dfrac{6}{8}$. Tổng hai tử số: $5 + 6 = 11$.\\
\textbf{Đáp án của câu hỏi là C}
}
\end{ex}

\begin{ex} %[4D10V57]
Quy đồng mẫu số của số 2 và phân số $\dfrac{3}{5}$ ta được
\choice
{\True $\dfrac{10}{5}$ và $\dfrac{3}{5}$}
{$\dfrac{5}{5}$ và $\dfrac{3}{5}$}
{$\dfrac{2}{5}$ và $\dfrac{3}{5}$}
{$\dfrac{6}{5}$ và $\dfrac{3}{5}$}
\loigiai{
$2 = \dfrac{2}{1} = \dfrac{10}{5}$.\\
\textbf{Đáp án của câu hỏi là A}
}
\end{ex}

\begin{ex} %[4D10V57]
Tìm mẫu số chung nhỏ nhất của hai phân số $\dfrac{5}{12}$ và $\dfrac{7}{18}$:
\choice
{$216$}
{$72$}
{$48$}
{\True $36$}
\loigiai{
$36$ chia hết cho cả 12 và 18.\\
\textbf{Đáp án của câu hỏi là D}
}
\end{ex}

\begin{ex} %[4D10C57]
Quy đồng mẫu số ba phân số $\dfrac{1}{3}, \dfrac{1}{4}, \dfrac{1}{5}$ được ba phân số có mẫu số chung nhỏ nhất là 60. Tổng các tử số mới là
\choice
{$45$}
{\True $47$}
{$49$}
{$51$}
\loigiai{
Ba phân số: $\dfrac{20}{60}, \dfrac{15}{60}, \dfrac{12}{60}$. Tổng tử số: $20 + 15 + 12 = 47$.\\
\textbf{Đáp án của câu hỏi là B}
}
\end{ex}

\textbf{Phần 2. Trắc nghiệm đúng sai}
\setcounterpageref{ex}{0}

\begin{ex} %[4D10H57]
Xét tính đúng sai khi quy đồng mẫu số hai phân số $\dfrac{3}{4}$ và $\dfrac{5}{6}$:
\choiceTF
{\True Có thể chọn mẫu số chung là $24$ ($4 \times 6$).}
{\True Có thể chọn mẫu số chung nhỏ nhất là $12$.}
{Khi quy đồng với MSC là $12$ ta được $\dfrac{6}{12}$ và $\dfrac{10}{12}$.}
{\True Khi quy đồng với MSC là $12$ ta được $\dfrac{9}{12}$ và $\dfrac{10}{12}$.}
\loigiai{
\begin{itemize}
    \item MSC 24: Đúng.
    \item MSC nhỏ nhất 12: Đúng.
    \item $3/4 = 9/12 \neq 6/12$. Ý 3 Sai, Ý 4 Đúng.
\end{itemize}
\begin{itemize}
    \item a) \textbf{Đúng}.
    \item b) \textbf{Đúng}.
    \item c) \textbf{Sai}.
    \item d) \textbf{Đúng}.
\end{itemize}
}
\end{ex}

\begin{ex} %[4D10V57]
Cho hai phân số $\dfrac{2}{5}$ và $\dfrac{4}{15}$. Xét tính đúng sai:
\choiceTF
{\True Mẫu số $15$ chia hết cho mẫu số $5$.}
{\True Chọn $15$ làm mẫu số chung.}
{Quy đồng được $\dfrac{8}{15}$ và $\dfrac{4}{15}$.}
{\True Quy đồng được $\dfrac{6}{15}$ và $\dfrac{4}{15}$.}
\loigiai{
\begin{itemize}
    \item 15 chia hết cho 5: Đúng.
    \item Chọn 15 làm MSC: Đúng.
    \item $2/5 = 6/15 \neq 8/15$. Ý 3 Sai, Ý 4 Đúng.
\end{itemize}
\begin{itemize}
    \item a) \textbf{Đúng}.
    \item b) \textbf{Đúng}.
    \item c) \textbf{Sai}.
    \item d) \textbf{Đúng}.
\end{itemize}
}
\end{ex}

\textbf{Phần 3. Trả lời ngắn}
\setcounterpageref{ex}{0}

\begin{ex} %[4D10H57]
Tìm mẫu số chung của hai phân số $\dfrac{2}{3}$ và $\dfrac{1}{7}$ (lấy tích hai mẫu số).
\shortans{21}
\loigiai{
\textbf{Đáp số:} $21$\\
$3 \times 7 = 21$.
}
\end{ex}

\begin{ex} %[4D10H57]
Quy đồng mẫu số phân số $\dfrac{3}{4}$ với mẫu số chung 20. Tìm tử số mới.
\shortans{15}
\loigiai{
\textbf{Đáp số:} $15$\\
$3 \times 5 = 15$.
}
\end{ex}

\begin{ex} %[4D10V57]
Quy đồng mẫu số hai phân số $\dfrac{5}{6}$ và $\dfrac{3}{8}$ với mẫu số chung nhỏ nhất là 24. Tìm tổng hai tử số mới.
\shortans{29}
\loigiai{
\textbf{Đáp số:} $29$\\
$5 \times 4 + 3 \times 3 = 20 + 9 = 29$.
}
\end{ex}

\begin{ex} %[4D10V57]
Tìm mẫu số chung nhỏ nhất của hai phân số $\dfrac{7}{12}$ và $\dfrac{5}{16}$.
\shortans{48}
\loigiai{
\textbf{Đáp số:} $48$\\
$48$.
}
\end{ex}

\begin{ex} %[4D10V57]
Quy đồng mẫu số của 3 và $\dfrac{2}{7}$ với mẫu số là 7. Tìm tử số của số 3 sau khi quy đồng.
\shortans{21}
\loigiai{
\textbf{Đáp số:} $21$\\
$3 \times 7 = 21$.
}
\end{ex}

\begin{ex} %[4D10C57]
Tìm mẫu số chung nhỏ nhất của 3 phân số: $\dfrac{1}{6}, \dfrac{1}{8}, \dfrac{1}{12}$.
\shortans{24}
\loigiai{
\textbf{Đáp số:} $24$\\
$24$.
}
\end{ex}
